% ════════════════════════════════════════════════════════════════
%  CHAPTER 2 : GEOMETRICAL PARAMETERS NL MECHANICS (2D)
% ════════════════════════════════════════════════════════════════
\chapter{Geometrically Parameterized Nonlinear Mechanics in 2D}
\label{ch:nl-mechanics-2d}

This chapter details the mathematical formulation and computational methodology for solving parametric, geometrically nonlinear elasticity problems in two dimensions using IGA and ROM.

% ──────────────────────────────────────────────────────────────
\section{Weak Formulation on a Reference Configuration}
\label{sec:weak-form}
% ──────────────────────────────────────────────────────────────

\subsection{Total Lagrangian Formulation}

All quantities are referred to the undeformed reference configuration $\Omega_0$. The principle of virtual work states: find $\vect{u} \in \mathcal{V}$ such that for all $\delta\vect{u} \in \mathcal{V}_0$:
\begin{equation}
    \int_{\Omega_0} \mat{S}(\vect{u}) : \delta\mat{E}(\vect{u}, \delta\vect{u}) \, \dd\Omega_0 = \int_{\Gamma_N} \vect{t} \cdot \delta\vect{u} \, \dd\Gamma + \int_{\Omega_0} \vect{b} \cdot \delta\vect{u} \, \dd\Omega_0
    \label{eq:weak-form}
\end{equation}
where $\mat{S}$ is the second Piola--Kirchhoff stress, $\delta\mat{E}$ is the variation of the Green--Lagrange strain, $\vect{t}$ represents prescribed tractions, and $\vect{b}$ body forces.

\subsection{Variation of the Green--Lagrange Strain}

The variation of $\mat{E}$ with respect to a virtual displacement $\delta\vect{u}$ is:
\begin{equation}
    \delta\mat{E} = \frac{1}{2}\left(\mat{F}^T \nabla\delta\vect{u} + (\nabla\delta\vect{u})^T \mat{F}\right)
\end{equation}

This expression is linear in $\delta\vect{u}$ but depends nonlinearly on the current displacement $\vect{u}$ through the deformation gradient $\mat{F}$.

\subsection{Parameterized Reference Domain}

When the geometry is parameterized by $\mu$ (e.g., aspect ratio, hole radius), the physical domain $\Omega(\mu)$ varies. In the IGA framework, this is handled naturally through parameter-dependent control points:
\begin{equation}
    \vect{x}(\xi, \mu) = \sum_{i=1}^{n} R_i(\xi) \, \vect{P}_i(\mu)
\end{equation}

All integrals are evaluated on the fixed master domain $\hat{\Omega} = [0,1]^2$, with the Jacobian $\mat{J}(\xi, \mu) = \partial\vect{x}/\partial\xi$ absorbing the parameter dependence:
\begin{equation}
    \int_{\Omega(\mu)} (\cdot) \, \dd\Omega = \int_{\hat{\Omega}} (\cdot) \, |\det\mat{J}(\xi, \mu)| \, \dd\hat{\Omega}
\end{equation}

% ──────────────────────────────────────────────────────────────
\section{IGA Discretization}
\label{sec:iga-discretization}
% ──────────────────────────────────────────────────────────────

\subsection{2D NURBS Approximation}

The displacement field is approximated using tensor-product NURBS:
\begin{equation}
    \vect{u}^h(\xi, \eta) = \sum_{A=1}^{N_{cp}} R_A(\xi, \eta) \, \vect{d}_A, \qquad R_A = \frac{N_i^p(\xi) \, M_j^q(\eta) \, w_{ij}}{W(\xi,\eta)}
\end{equation}
where $N_{cp} = n \times m$ is the total number of control points and $W(\xi,\eta)$ is the weighting function.

\subsection{Stiffness Matrix Assembly}

The tangent stiffness matrix is assembled element-by-element using Gauss quadrature on the master domain:
\begin{equation}
    K_{AB}^{IJ} = \int_{\hat{\Omega}} \pd{R_A}{x_I} \, C_{IJKL} \, \pd{R_B}{x_K} \, |\det\mat{J}| \, \dd\hat{\Omega}
\end{equation}

Physical derivatives are obtained through the Jacobian inverse:
\begin{equation}
    \begin{bmatrix} \partial R_A / \partial x \\ \partial R_A / \partial y \end{bmatrix} = \mat{J}^{-1} \begin{bmatrix} \partial R_A / \partial \xi \\ \partial R_A / \partial \eta \end{bmatrix}
\end{equation}

For plane stress conditions, the elasticity matrix is:
\begin{equation}
    \mat{C} = \frac{E}{1 - \nu^2} \begin{bmatrix} 1 & \nu & 0 \\ \nu & 1 & 0 \\ 0 & 0 & \frac{1-\nu}{2} \end{bmatrix}
\end{equation}

\subsection{Boundary Conditions}

Dirichlet boundary conditions are applied using a penalty method with factor $\beta \approx 10^{12}$:
\begin{equation}
    \Pi_{\text{penalty}} = \frac{1}{2} \beta \, |u_i - \bar{u}_i|^2
\end{equation}

This approach is also used to enforce compatibility at degenerate NURBS points (e.g., collapsed corners in the plate-with-hole geometry).

% ──────────────────────────────────────────────────────────────
\section{Offline Simulations}
\label{sec:offline}
% ──────────────────────────────────────────────────────────────

\subsection{Parameter Space Definition}

The parameter vector $\mu \in \mathcal{P}$ is defined as:

\begin{table}[H]
\centering
\caption{Parameter space definition for the 2D test cases.}
\label{tab:param-space}
\begin{tabular}{llll}
\toprule
\textbf{Parameter} & \textbf{Symbol} & \textbf{Description} & \textbf{Range} \\
\midrule
Young's Modulus & $E$ & Material stiffness & $[10^6, 2 \times 10^8]$\,Pa \\
Poisson's Ratio & $\nu$ & Transverse strain ratio & $[0.25, 0.49]$ \\
Load Intensity & $F$ & Applied force magnitude & $[0, 1000]$\,N \\
Aspect Ratio & $\alpha = L/H$ & Geometry parameter & $[5, 20]$ \\
\bottomrule
\end{tabular}
\end{table}

\subsection{Latin Hypercube Sampling}

The training set $\{\mu_1, \ldots, \mu_{N_s}\}$ is generated using \acrfull{lhs} to ensure uniform coverage of the parameter space. An initial set of 25 snapshots is used for POD validation, with 100+ snapshots for the final hyper-reduced model.

\subsection{Snapshot Database Generation}

For each parameter sample $\mu_k$, the full-order model is solved using Newton--Raphson iteration:
\begin{equation}
    \mat{S} = \left[\vect{u}(\mu_1), \, \vect{u}(\mu_2), \, \ldots, \, \vect{u}(\mu_{N_s})\right] \in \R^{2N_{cp} \times N_s}
\end{equation}

Each column contains the converged displacement field (both $u_x$ and $u_y$ components interleaved) for one parameter configuration.

\subsection{POD Basis Extraction}

The SVD of the snapshot matrix yields the reduced basis:
\begin{equation}
    \mat{S} = \mat{U} \mat{\Sigma} \mat{V}^T, \qquad \mat{\Phi} = \mat{U}_{[:, 1:r]}
\end{equation}

The number of retained modes $r$ is determined by the energy criterion (Eq.~\ref{eq:energy-criterion}). For the 2D plate benchmarks, typically $r \leq 5$ modes capture $>99.99\%$ of the system energy.

% ──────────────────────────────────────────────────────────────
\section{Online Simulations}
\label{sec:online}
% ──────────────────────────────────────────────────────────────

\subsection{Galerkin-Projected Nonlinear ROM}

For a new parameter $\mu^*$, the online phase solves the reduced nonlinear system iteratively:

\begin{algorithm}[H]
\caption{Online NL-ROM Query}
\label{alg:nl-rom}
\begin{algorithmic}[1]
\REQUIRE Reduced basis $\mat{\Phi}$, parameter $\mu^*$, tolerance $\varepsilon$
\STATE Initialize $\vect{q}^{(0)} = \vect{0}$
\FOR{$k = 0, 1, 2, \ldots$}
    \STATE Reconstruct: $\vect{u}^{(k)} = \mat{\Phi} \, \vect{q}^{(k)}$
    \STATE Assemble $\mat{K}_T(\vect{u}^{(k)}, \mu^*)$ and $\vect{F}_{\text{int}}(\vect{u}^{(k)}, \mu^*)$
    \STATE Project: $\mat{K}_{\text{red}} = \mat{\Phi}^T \mat{K}_T \mat{\Phi}$, \quad $\vect{R}_{\text{red}} = \mat{\Phi}^T (\vect{F}_{\text{ext}} - \vect{F}_{\text{int}})$
    \STATE Solve: $\mat{K}_{\text{red}} \, \Delta\vect{q} = \vect{R}_{\text{red}}$
    \STATE Update: $\vect{q}^{(k+1)} = \vect{q}^{(k)} + \Delta\vect{q}$
    \IF{$\norm{\Delta\vect{q}} < \varepsilon$}
        \STATE \textbf{break} (converged)
    \ENDIF
\ENDFOR
\RETURN $\vect{u} = \mat{\Phi} \, \vect{q}$
\end{algorithmic}
\end{algorithm}

\subsection{Performance Targets}

\begin{table}[H]
\centering
\caption{Performance comparison: FOM vs.\ ROM.}
\label{tab:performance}
\begin{tabular}{lcc}
\toprule
\textbf{Metric} & \textbf{FOM} & \textbf{ROM} \\
\midrule
System size & $2000 \times 2000$ & $10 \times 10$ \\
Solve time per query & $\sim 2$\,s & $< 10$\,ms \\
Memory footprint & $\sim 50$\,MB & $\sim 10$\,KB \\
Relative $L^2$ error & --- & $< 1\%$ \\
\bottomrule
\end{tabular}
\end{table}

\subsection{Data Serialization}

Offline results (reduced basis $\mat{\Phi}$, projected operators, mesh metadata) are serialized to portable JSON format for consumption by the web-based Digital Twin frontend.
